\SourceLeaf{287}
\hspace*{\parindent}Die Logarithmen der Seiten des Dreiecks in Toisen sind
\begin{center}
\begin{tabular}{lr}
Hohehagen-Inselsberg&\(4{,}6393865\)\\
Inselsberg-Brocken&\(4{,}7353929\)\\
Brocken-Hohehagen&\(4{,}5502669\)
\end{tabular}
\end{center}
\hspace*{\parindent}Die Azimuthe sind
\begin{center}
\begin{tabular}{lrrr}
&\multicolumn{3}{c}{Standpunkt Brocken}\\
Inselsberg&\(5^0\)&\(42'\)&\(22''\)\\
Hohehagen&58&49&8\\
&\multicolumn{3}{c}{Standpunkt Hohehagen}\\
Brocken&238&9&2\\
Inselsberg&324&23&1\\
&\multicolumn{3}{c}{Standpunkt Inselsberg}\\
Hohehagen&144&55&51\\
Brocken&185&35&21
\end{tabular}
\end{center}
\hspace*{\parindent}Man braucht hiebei zwischen Werthen auf dem Sphaeroid und denen auf der Kugel nicht zu unterscheiden, da für die Logarithmen der Abstände erst in der achten oder neunten Decimale, für die Azimuthe erst in dem Tausendtheilen der Secunde Ungleichheit eintritt, und für unsern Zweck Logarithmen mit vier Decimalen und Azimuthe in Minuten schon überflüssig genau sind. Die Rechnung nach obigen Formeln gibt hiermit folgende Reductionen, wie sie mit ihren Zeichen zu den Azimuthen auf dem Sphaeroid addirt werden müssen, um die Azimuthe auf der Kugel zu erhalten:
\begin{center}
\begin{tabular}{lr}
Brocken-Inselsberg&\(+0''00055\)\\
Brocken-Hohehagen&\(+0{,}00196\)\\
Hohehagen-Brocken&\(-0{,}00238\)\\
Hohehagen-Inselsberg&\(-0{,}00332\)\\
Inselsberg-Hohehagen&\(+0{,}00428\)\\
Inselsberg-Brocken&\(-0{,}00083\)
\end{tabular}
\end{center}
\hspace*{\parindent}Die Winkel des Dreiecks auf dem Sphaeroid (zwischen den geodätischen Linien) empfangen also zur Reduction auf die Winkel des Kugeldreiecks (zwischen Grössten Kreisbögen) die Aenderungen
